% Standalone-Variablendiagramm (Forschungsdesign).
% Kompilieren:  pdflatex variablendiagramm.tex
% PNG:          pdftoppm -png -r 300 variablendiagramm.pdf variablendiagramm
\documentclass[border=12pt]{standalone}
\usepackage[utf8]{inputenc}
\usepackage[T1]{fontenc}
\usepackage[ngerman]{babel}
\usepackage{tikz}
\usetikzlibrary{positioning,arrows.meta,calc}

\definecolor{cData}{HTML}{DCE7F5}
\definecolor{cDataB}{HTML}{5B7FB0}
\definecolor{cEval}{HTML}{EADCF2}
\definecolor{cEvalB}{HTML}{8A5FA8}
\definecolor{cCtrl}{HTML}{ECECEC}
\definecolor{cCtrlB}{HTML}{9A9A9A}
\definecolor{cLbl}{HTML}{6E6E6E}

\begin{document}
\begin{tikzpicture}[
  font=\small\sffamily,
  box/.style={draw, rounded corners=2pt, align=center, minimum height=11mm,
              inner sep=6pt, line width=0.6pt},
  uv/.style={box, fill=cData, draw=cDataB, text width=44mm},
  av/.style={box, fill=cEval, draw=cEvalB, text width=40mm, minimum height=16mm},
  ctrl/.style={box, fill=cCtrl, draw=cCtrlB, text width=92mm},
  lbl/.style={font=\footnotesize\itshape\sffamily, text=cLbl},
  arr/.style={-{Stealth[length=2.8mm]}, line width=0.8pt},
  arrc/.style={-{Stealth[length=2.6mm]}, line width=0.7pt, dashed, draw=cCtrlB},
]

% --- Unabhaengige Variablen -------------------------------------------------
\node[uv] (uva) {\textbf{UV$_1$: Datenquelle}\\ 5 Stufen\\ \footnotesize (emobpy, real\_world\_ev, ved, routine, yjmob100k)};
\node[uv, below=16mm of uva] (uvb) {\textbf{UV$_2$: Algorithmus}\\ 3 Stufen\\ \footnotesize (Random Forest, LightGBM, LSTM)};

% --- Abhaengige Variable (rechts, mittig) -----------------------------------
\coordinate (mid) at ($(uva)!0.5!(uvb)$);
\node[av, right=42mm of mid] (av) {\textbf{AV: Prognosegüte}\\ \footnotesize Klassifikationsmetriken,\\ \footnotesize P10/P90-Quantile};

% --- Kontrollvariablen ------------------------------------------------------
\node[ctrl, below=14mm of uvb.south west, anchor=north west] (ctrl)
  {\textbf{Kontrollvariablen} (fixiert)\\ Hyperparameter \textperiodcentered\ Feature-Set \textperiodcentered\ Trainingsprotokoll};

% --- Kanten -----------------------------------------------------------------
\draw[arr] (uva.east) -- (av.west) node[pos=0.55, above, sloped, lbl] {variiert};
\draw[arr] (uvb.east) -- (av.west) node[pos=0.55, below, sloped, lbl] {variiert};
\draw[arrc] (ctrl.north) -- ++(0,20pt) -| (av.south) node[pos=0.4, above, lbl] {konstant gehalten};

\end{tikzpicture}
\end{document}
